% coverpage_example.tex - Test for Cover Page Enhancements (v7.1.0)
% This example demonstrates all v7.1.0 cover page features
% Compiler: LuaLaTeX
%
% Features demonstrated:
%   - Decorative double frame on title page
%   - Cover image with faded edges
%   - Hebrew signature in handwriting font (Guttman Yad)
%   - Signature image below handwriting name
%   - \makesignature command for signature block

\documentclass[book]{hebrew-academic-template}

% Graphics path for images
\graphicspath{{images/}}

%% ============================================
%% Book Metadata (v7.1.0 commands)
%% ============================================
\hebrewtitle{דוגמה לעמוד שער משופר}
\englishtitle{Enhanced Cover Page Example}
\hebrewsubtitle{הדגמת תכונות גרסה \en{7.1.0}}
\hebrewauthor{ד"ר יורם סגל}
\hebrewversion{גרסה \textenglish{1.0}}
\date{\textenglish{January 2026}}

% v7.1.0: Cover image with faded edges (path relative to graphicspath)
\coverimage{fig_architecture.png}

% v7.1.0: Hebrew signature (uses Guttman Yad font if available)
\hebrewsignature{יורם}

% v7.1.0: Signature image
\signatureimage{signature_example.png}

\begin{document}

%% ============================================
%% Front Matter
%% ============================================
\frontmatter

% Title Page - automatically includes:
%   - Decorative double frame (outer 3pt + inner 1pt)
%   - Cover image with faded edges (if \coverimage set)
\maketitle

%% ============================================
%% Personal Word with Signature
%% ============================================
\newpage
\thispagestyle{empty}

{\LARGE\bfseries מילה אישית}
\addcontentsline{toc}{chapter}{מילה אישית}

\bigskip

אני שמח להציג בפניכם את גרסה \en{7.1.0} של התבנית האקדמית העברית. גרסה זו כוללת שיפורים משמעותיים בעמוד השער:

\begin{itemize}
    \item \textbf{מסגרת כפולה מעוטרת} -- מסגרת חיצונית עבה (\en{3pt}) ופנימית דקה (\en{1pt})
    \item \textbf{תמונת שער עם גבולות דהויים} -- אפקט \en{fading} על ארבעת הצדדים
    \item \textbf{חתימה בכתב יד} -- שימוש בגופן \en{Guttman Yad} לשם בעברית
    \item \textbf{תמונת חתימה} -- אפשרות להוסיף תמונת חתימה סרוקה
\end{itemize}

\vspace{1cm}

% v7.1.0: Signature block - renders Hebrew name in handwriting font
% Aligned left with 5em offset from margin
\makesignature

%% ============================================
%% Table of Contents
%% ============================================
\tableofcontents

%% ============================================
%% Main Matter
%% ============================================
\mainmatter

\hebrewchapter{מבוא לתכונות החדשות}

פרק זה מתאר את השימוש בתכונות עמוד השער החדשות בגרסה \en{7.1.0}.

\hebrewsection{פקודות מטא-דאטה}

הפקודות הבאות זמינות בפריאמבל:

\begin{itemize}
    \item \en{\textbackslash hebrewtitle\{...\}} -- כותרת עברית (חובה)
    \item \en{\textbackslash englishtitle\{...\}} -- כותרת אנגלית (חובה)
    \item \en{\textbackslash hebrewsubtitle\{...\}} -- כותרת משנה (אופציונלי)
    \item \en{\textbackslash hebrewauthor\{...\}} -- שם המחבר (אופציונלי)
    \item \en{\textbackslash hebrewversion\{...\}} -- גרסה (אופציונלי)
    \item \en{\textbackslash coverimage\{path\}} -- תמונת שער (אופציונלי, \en{v7.1.0})
    \item \en{\textbackslash hebrewsignature\{name\}} -- שם בכתב יד (אופציונלי, \en{v7.1.0})
    \item \en{\textbackslash signatureimage\{path\}} -- תמונת חתימה (אופציונלי, \en{v7.1.0})
\end{itemize}

\hebrewsection{פקודות רינדור}

\begin{itemize}
    \item \en{\textbackslash maketitle} -- מרנדר את עמוד השער עם המסגרת
    \item \en{\textbackslash makesignature} -- מרנדר את בלוק החתימה
\end{itemize}

\hebrewsection{דוגמת קוד}

\begin{english}
\begin{verbatim}
% In preamble:
\hebrewtitle{כותרת בעברית}
\englishtitle{English Title}
\coverimage{images/cover.png}
\hebrewsignature{יורם}
\signatureimage{images/signature.png}

% In document:
\maketitle              % Renders title page with frame
...
\makesignature          % Renders signature block
\end{verbatim}
\end{english}

\hebrewchapter{גופן כתב יד}

\hebrewsection{גופן \en{Guttman Yad}}

הפקודה \en{\textbackslash hebrewsignature} משתמשת בגופן \en{Guttman Yad} אם הוא מותקן במערכת.

אם הגופן לא מותקן, הטקסט יוצג בגופן \textit{\textbf{נטוי מודגש}} כחלופה.

\hebrewsection{התקנת הגופן}

גופן \en{Guttman Yad} הוא גופן עברי חינמי שניתן להוריד מהאינטרנט. לאחר ההתקנה, הוא יזוהה אוטומטית.

%% ============================================
%% Back Matter
%% ============================================
\backmatter

\end{document}
